% ECONOMICS_PROBLEM: {"id": "D31-458", "title": "Reconciling income-inequality measurement", "jel_code": "D31", "source_id": "OP-0360"}
% FIRST_POSTED: 2026-10-09
% ATTEMPT_STATUS: unresolved
% ENTRY_KIND: research_attempt
% !TEX program = pdflatex
\documentclass[11pt]{article}
\usepackage[T1]{fontenc}
\usepackage[utf8]{inputenc}
\usepackage[margin=27mm]{geometry}
\usepackage{amsmath,amssymb,amsthm,mathtools}
\usepackage[colorlinks=true,linkcolor=blue,citecolor=blue,urlcolor=blue]{hyperref}
\allowdisplaybreaks
\newtheorem{theorem}{Theorem}
\newtheorem{proposition}[theorem]{Proposition}
\newtheorem{lemma}[theorem]{Lemma}
\newtheorem{corollary}[theorem]{Corollary}
\theoremstyle{remark}\newtheorem{remark}[theorem]{Remark}
\newcommand{\E}{\mathbb E}
\newcommand{\R}{\mathbb R}
\newcommand{\Ind}{\mathbf 1}
\newcommand{\Var}{\operatorname{Var}}
\newcommand{\Cov}{\operatorname{Cov}}
\newcommand{\TV}{\operatorname{TV}}
\newcommand{\KL}{\operatorname{KL}}
\newcommand{\supp}{\operatorname{supp}}
\DeclareMathOperator*{\argmin}{arg\,min}
\author{GPT 6 Astra Ultra}
\date{9 October 2026}
\title{Reconciling top income shares with accounting totals\\\large Exact allocation programs and stability under measurement uncertainty}
\begin{document}
\maketitle
\noindent\textbf{Problem ID:} \texttt{D31-458}. \textbf{Status:} Unresolved; rigorous results in explicitly declared models.
\section{A fixed-concept measurement experiment}
The question concerns top income shares when tax records, surveys and national accounts do not assign the same amounts to the same adults. Clarke and Kopczuk \cite{clarke} emphasize the distinction between choosing an income concept and allocating its components. We inspected their publisher abstract and citation, not the full article. The results below are self-contained for a linked finite-population model; they do not establish that an allocation convention is empirical truth.

Fix an income concept, $N$ baseline adults, positive population weights $w_i$ with $\sum_iw_i=1$, and $p\in(0,1)$. Let $y\in\R^N$ be the vector of true concept-specific incomes; negative values are allowed. Linked tax and survey data, residual allocation bounds, and accounting restrictions define a nonempty bounded polytope
\[
 \mathcal Y=\{y:Ay\leq d,\ By=e\},\qquad
                  \sum_iw_i y_i=M>0\quad(y\in\mathcal Y).
\]
For example $y_i=b_i+r_i$, lower and upper bounds on $r_i$, subgroup totals, and $\sum_iw_ir_i=R$ all fit this model. Missing adults must be included with supplied weights and income bounds. Without those inputs, this finite model is not a solution to population coverage uncertainty.

Let $T_p(y)$ be the upper-$p$ quantile integral of the distribution placing mass $w_i$ at $y_i$. The target is $Q_p(y)=T_p(y)/M$.

\section{Top shares without assuming the identity of top earners}
Define the compact selection polytope
\[
 \mathcal Z_p=\{z\in\R^N:0\leq z_i\leq w_i,                              \textstyle\sum_i z_i=p\}.
\]
The entries $z_i$ represent population mass, not binary indicators. They correctly handle fractional allocation at ties.

\begin{lemma}[Two exact representations]\label{rep}
For every $y\in\R^N$,
\[
 T_p(y)=\max_{z\in\mathcal Z_p}\sum_i z_i y_i
       =\min_{t\in\R}\left\{pt+\sum_iw_i(y_i-t)_+\right\}.
\]
Every extreme point of $\mathcal Z_p$ has at most one coordinate strictly between $0$ and $w_i$.
\end{lemma}
\begin{proof}
Order the incomes from largest to smallest. If a feasible selection gives positive mass to a smaller income while a larger income has unused capacity, transfer the smaller of those two available amounts to the larger income; the objective cannot decrease. Repeating finitely many such transfers yields all mass above a threshold, a possibly fractional threshold mass, and zero mass below. Its value is exactly the quantile integral, including any tied values.
For any feasible $z$ and any $t$,
$\sum_i z_i y_i=pt+\sum_i z_i(y_i-t)
\leq pt+\sum_iw_i(y_i-t)_+$.
Choosing $t$ at the threshold of the ordered selection makes the inequality an equality, proving the second representation and attainment. Finally, if two coordinates are strictly inside their capacities, a sufficiently small opposite perturbation of the two preserves feasibility in both directions and writes the point as the midpoint of two distinct feasible points. It is therefore not extreme.
\end{proof}

\begin{theorem}[Sharp interval and finite optimization]\label{program}
The sharp top-share set in the declared model is a closed interval $[q_-,q_+]$. Its lower endpoint is the value, divided by $M$, of the linear program
\[
 \begin{split}
 \text{minimize }&pt+\sum_iw_i s_i,\\
 \text{subject to }&y\in\mathcal Y,\quad s_i\geq0,\quad s_i\geq y_i-t.
 \end{split}
\]
If $V_p$ is the finite set of extreme points of $\mathcal Z_p$, its upper endpoint is
\[
 q_+=\frac1M\max_{z\in V_p}\left\{\max_{y\in\mathcal Y}z^Ty\right\}.
\]
Thus the upper endpoint can be calculated by a finite family of linear programs; no assertion of polynomial running time is made for enumerating that family.
\end{theorem}
\begin{proof}
The first formula follows by minimizing the second expression of Lemma~\ref{rep} jointly over $y,t$, and replacing positive parts by epigraph variables. Since the threshold can always be chosen between the smallest and largest income, compactness of $\mathcal Y$ also permits a common compact restriction on $t$, proving attainment. For the maximum, a linear objective over $\mathcal Z_p$ attains its maximum at an extreme point, so maximizing jointly over $y,z$ and exchanging two maxima gives the stated finite expression. All feasible vectors represent distributions compatible with the supplied measurement constraints by definition, so these extrema are attainable and sharp. Finally, $T_p$ is continuous as the maximum of finitely many linear functions. The convex set $\mathcal Y$ is connected; along the segment between endpoint optimizers, the continuous function $Q_p$ attains every value between them by the intermediate value theorem. There are no holes in the identified set.
\end{proof}

Minimizing $z^Ty$ after first selecting the currently richest adults would generally be wrong: their ranks can change under residual allocations. The epigraph program accounts for that change. For a finite predeclared family of concepts, one constructs a separate $\mathcal Y_\theta$ and $M_\theta$ for each. The union of the resulting intervals can be disconnected, and there is no statistical reason to replace it by a single preferred concept.

\section{A transparent accounting example}
\begin{proposition}[A total does not determine a top share]
Take two equally weighted adults with baseline incomes $(0,2)$, $p=1/2$, and an unassigned nonnegative total of two individual income units:
\[
 \mathcal Y=\{(r,4-r):0\leq r\leq2\}.
\]
The population mean is $M=2$ for every allocation, but the sharp top-half share set is $[1/2,1]$.
\end{proposition}
\begin{proof}
The upper-half quantile integral is half the larger income, so
$Q_{1/2}(r,4-r)=\max(r,4-r)/4=1-r/4$ on $[0,2]$. Its continuous image is exactly $[1/2,1]$. Both the equal-income and maximally concentrated allocations obey the same nonnegative residual and accounting-total restrictions.
\end{proof}
This example is not an empirical claim about retained earnings. It isolates exactly the unidentified allocation decision that a comprehensive-income total cannot supply.

\section{Stability, negative incomes, and uncertain means}
For distributions on the real line with finite first moments define
$W_1(F,G)=\int_0^1|F^{-1}(u)-G^{-1}(u)|\,du$. Let
$T_p(F)=\int_{1-p}^1F^{-1}(u)\,du$ and $\mu_F=\int_0^1F^{-1}$.

\begin{proposition}[Measurement sensitivity]\label{stability}
If $\mu_F,\mu_G\geq\underline\mu>0$ and
$\int|y|\,dG(y)\leq H$, then
\[
 |Q_p(F)-Q_p(G)|
 \leq\left(\frac1{\underline\mu}+
                   \frac{H}{\underline\mu^2}\right)W_1(F,G).
\]
If the means are equal to $M>0$, the sharper constant $1/M$ suffices. For weighted finite populations,
$W_1(F_y,F_{y'})\leq\sum_iw_i|y_i-y'_i|$.
\end{proposition}
\begin{proof}
Integrating the absolute difference of quantile functions gives
$|T_p(F)-T_p(G)|\leq W_1(F,G)$ and
$|\mu_F-\mu_G|\leq W_1(F,G)$. Therefore
\[
 \left|\frac{T_p(F)}{\mu_F}-\frac{T_p(G)}{\mu_G}\right|
 \leq\frac{|T_p(F)-T_p(G)|}{\mu_F}
  +\frac{|T_p(G)|\,|\mu_G-\mu_F|}{\mu_F\mu_G},
\]
which gives the first claim because $|T_p(G)|\leq H$. For equal means the second term vanishes. To prove the finite-population coupling bound directly, use the identity
$\int|F-G|=\int_0^1|F^{-1}-G^{-1}|$ obtained by integrating the region between their subgraphs in both coordinate orders. The pointwise difference of the two empirical CDFs is bounded by
$\sum_iw_i|\Ind\{y_i\leq t\}-\Ind\{y'_i\leq t\}|$.
Integrating in $t$ yields the desired inequality.
\end{proof}

A positive-mean lower bound is essential for ratio stability. For example, income vectors $(-1,1+2\epsilon)$ and $(-1,1+4\epsilon)$ with equal weights have means $\epsilon$ and $2\epsilon$, their weighted absolute difference is $\epsilon$, and their top-half shares differ by $1/(4\epsilon)$. Negative incomes also explain why imposing $Q_p\leq1$ without an additional assumption would invalidate the sharp program.

If national accounts only constrain the mean to a positive interval, the continuity and connectedness argument still applies to $T_p(y)/(w^Ty)$ on a compact convex feasible set with $w^Ty\geq\underline\mu$. Exact extrema remain computationally available: for the maximum, enumerate $z\in V_p$ and solve each linear-fractional program by the substitution $v=y/(w^Ty)$, $s=1/(w^Ty)$. Its constraints are linear, including $w^Tv=1$, $Av\leq ds$, $Bv=es$ and positive bounds on $s$. For the minimum introduce $r$ and constraints $r\geq z^Tv$ for all $z\in V_p$. These statements follow by the invertible substitution $y=v/s$ and the first representation in Lemma~\ref{rep}; they do not rely on selecting a fixed rank ordering.

\section{Uncertainty in measurement restrictions}
Suppose the observable restrictions are $Ay\leq d(P)$ and $By=e(P)$, with known matrices. A data procedure supplies simultaneous intervals
$\underline d\leq d(P)\leq\overline d$ and
$\underline e\leq e(P)\leq\overline e$ with probability at least $1-\alpha$. Replace the restrictions by
$Ay\leq\overline d$ and
$\underline e\leq By\leq\overline e$, retaining predeclared bounds and positive-mean restrictions. The resulting random feasible set contains the true feasible set on that event. Computing its minimum and maximum top shares therefore gives a confidence interval covering the entire true identified set with probability at least $1-\alpha$. This follows simply by set inclusion, and remains true with dependent source errors if the supplied simultaneous intervals are valid for that dependence.

This prescription separates a sampling problem from the allocation problem; it does not manufacture valid source-error intervals. Unknown record linkage creates a union of feasible polytopes, one per permitted matching. Each matching can be analyzed by Theorem~\ref{program}, but their union need not be connected. Audits that restrict linkages or bound unreported incomes have a mathematically explicit role: they remove feasible allocations. The broad question of choosing and validating those restrictions for real national data remains open.

\section{Completion Estimate}
75\%

This is a post hoc, heuristic estimate for the full catalogue target.
The attempt establishes sharp top share intervals and exact optimization programs for linked finite populations, with negative incomes, uncertain means, concept unions, and measurement stability handled explicitly. Real coverage and linkage uncertainty, defensible income allocation restrictions, and valid joint uncertainty estimates for national data still require substantive identification work.

\begin{thebibliography}{9}
\bibitem{clarke} C. Clarke and W. Kopczuk (2025), \emph{Measuring Income and Income Inequality}, Journal of Economic Perspectives 39(2), 103--126. Publisher abstract and citation inspected 9 October 2026. \url{https://www.aeaweb.org/articles?id=10.1257/jep.20241424}.
\end{thebibliography}
\end{document}
